% ============================================================================
% Section 3: Methods (Phase B, written from docs/PAPER_OUTLINE.md section 3).
% Every number is a generated macro (experiments/FIGS/make_tables.py); the
% class counts of the label sets are \Kfull, \Kres, \Kmid, \Kcoarse.
% ============================================================================
\section{Methods}\label{sec:methods}

\subsection{Data, model and pretraining}\label{sec:data}

We pretrain on JetClass-II~\cite{sophon,jetclasstwo}, the dataset of the published Sophon model, with its
selection and its reweighting, so that every label set learns from the jet distribution Sophon learned from.
JetClass-II holds about \LitJetClassTwoNJetsMillion{} million simulated jets, clustered with the
anti-$k_{\mathrm T}$ algorithm~\cite{antikt} with $R = \LitJetClassTwoJetRadius$. They are labeled with the
\LitJetClassTwoNResonantClasses{} resonant and \LitJetClassTwoNQcdClasses{} QCD classes of JetClass-II. In a
resonant jet a spin-0 resonance $X$ decays either directly into a pair of quarks, gluons or leptons, or into a
pair of intermediate resonances $Y$ that decay in turn. The masses of $X$, up to
\LitJetClassTwoMaxResonanceMass~GeV, and of $Y$ are sampled for each event, so a class is defined by its decay
products and not by a mass. Resonances are generated with MadGraph5\_aMC@NLO~\cite{madgraph} and QCD multijet
events with Pythia~\LitPythiaVersion~\cite{pythia}, which also simulates the decays, the parton shower and
hadronization. Delphes~\cite{delphes} simulates a CMS-like detector with an average of
\LitJetClassTwoPileupMean{} additional proton--proton interactions per event, whose particles
PUPPI~\cite{puppi} suppresses.

We keep jets with $\DesignPtMin < p_{\mathrm T} < \DesignPtMax$~GeV and $\DesignMsdMin < m_{\mathrm{SD}} <
\DesignMsdMax$~GeV. Here $m_{\mathrm{SD}}$ is the soft-drop mass, the jet mass after soft, wide-angle
radiation is removed, which peaks near the mass of the decaying particle. It is computed with the soft-drop
procedure~\cite{softdrop} at $z_{\mathrm{cut}} = \LitCmsSdZcut$ and $\beta = \LitCmsSdBeta$, the values CMS
uses~\cite{cmsmlid}. Every JetClass-II jet also satisfies $|\eta| < \LitSophonEtaMax$~\cite[App.~A]{sophon}.
These requirements leave \DesignNTrainSelected{} training jets. % CHECK: pending DesignNTrainSelected (ground truth 110,087,593; no committed count yet)
Training jets are drawn with Sophon's reweighting, which gives each of Sophon's reweighting categories the
same distribution in $(p_{\mathrm T}, m_{\mathrm{SD}})$ in order to avoid sculpting the background mass
distribution~\cite[Sec.~II]{sophon}. The categories are fixed by the JetClass-II classes and not by the label
set, so which jets are drawn, and how often, is the same for every label set. Every linear probe and every
fine-tuning on JetClass-II uses jets from files that no pretraining run read.

The network is Sophon's~\cite{sophon}, built on the Particle Transformer (ParT) architecture~\cite{part}. It has
\DesignParticleBlocks{} particle-attention blocks, whose attention carries a bias computed from pairwise
features of the particles, \DesignClassBlocks{} class-attention blocks, an embedding dimension of
\DesignEmbedDim, and a classification head with one hidden layer of \DesignHeadHidden{} ReLU units. The number
of particle-attention blocks follows the released code~\cite{sophoncode}: the Sophon paper states
\LitSophonPaperParticleBlocks{} blocks, but its \LitSophonParamsMillion{} million parameters correspond to
\LitSophonBlocksCode. As in Sophon, the four-momenta of the particles are rescaled so that every jet has the
same transverse momentum. The frozen features, the output of the backbone, are the \DesignEmbedDim-dimensional
class token after the class-attention blocks and the final layer normalization, which is the input of the head. The initial weights
follow PyTorch's default initialization rather than the one in Sophon's code. Absolute performance therefore
need not match Sophon's model, while every comparison in this paper is between models initialized the same
way.

Every configuration is pretrained with the same settings, so that within a seed the label set is the only
difference between two models (figure~\ref{fig:design}(c)). Models with the same seed share their initial
weights and see the same training examples in the same order. The order of the training jets depends only on
the seed and the epoch, and every layer starts from the same weights except the output layer, whose width is
the number of classes. Each model sees \DesignExamplesSeen{} training examples, \DesignEpochs{} epochs of
\DesignSamplesPerEpoch{} jets as in Sophon's released configuration~\cite{sophoncode}. It is trained with the
Ranger optimizer~\cite{lookahead,radam} and a learning rate that is constant at first and then decays
exponentially; appendix~\ref{sec:app-pretrain} gives the values. The checkpoint rule is shared too: each model is
evaluated at the first maximum of its validation accuracy within epochs
\DesignWavgFirstEpoch--\DesignWavgLastEpoch, where the learning rate is lowest. For every run, the code release
records a hash of the jet identifiers read in each epoch, against which the training order of any two models
can be checked.

% ----------------------------------------------------------------------------
\subsection{Label sets}\label{sec:labelsets}

The label set is the one variable of the study. The four label sets we pretrain form the class hierarchy of
figure~\ref{fig:design}(a), and each merge removes a named set of physics distinctions. The \Kfull-class label set is Sophon's. It has one class for each resonant decay
and divides the QCD jets into \LitJetClassTwoNQcdClasses{} classes by their numbers of $b$, $c$ and $s$
quarks. The \Kres-class level merges these QCD classes into one and keeps every resonant class, so it removes
only the quark content of QCD jets. The \Kmid-class level labels each decay by its topology, the pattern of
intermediate resonances and prongs, and by whether it contains a $b$ quark, else a $c$ quark, else only light
quarks or gluons. It removes three kinds of distinction at once: the number of $b$ or $c$ quarks, gluon
against light quark, and electron against muon, also for $\tau$ leptons that decay to either. Grouping gluons
with light quarks is our choice; Sophon's classes keep them apart. The \Kcoarse-class level removes quark
flavor. It keeps, for each decay, its number of prongs and the type of each: quark or gluon, electron or muon,
$\tau$ lepton decaying to hadrons or to a lepton, or neutrino.

The levels remove flavor before structure, so the coarsest still describes the decay topology an analysis
selects on. At every level a class holds only QCD jets, only decays into quarks and gluons, or only decays with
leptons. The step from \Kres{} to \Kmid{} classes is the only one that removes more than one kind of
distinction, so no result in this paper attributes an effect of that step to one of its three.

Merging classes removes analysis discriminants outright, not only information. Analyses build discriminants by
summing the output scores of classes, and for an ideal classifier the score of a merged class is the sum of
the scores of the classes it merges~\cite[Eq.~(2)]{sophon}. A discriminant can therefore be formed exactly from
a label set only when every set of classes it sums is a union of that label set's classes, and
table~\ref{tab:usecase} lists the outcome for published discriminants. Sophon's $X\to b\bar b$
discriminant~\cite[Eq.~(4)]{sophon} can be formed at \Kfull{} and \Kres{} classes but not at \Kmid, where
$X\to b\bar b$ shares a class with $X\to bc$ and $X\to bq$. Sophon's other
discriminants~\cite[Eqs.~(5)--(9)]{sophon} exist only at \Kfull{} and \Kres{} classes. So does the
discriminant $D_{bc}$ of the boosted $|V_{cb}|$ measurement, whose denominator contains $X\to bq$ and
$X\to cs$~\cite[Eq.~(1)]{vcbboosted}. A discriminant that weights each QCD subclass separately, as GN3X
evaluates its rejection for each QCD subclass~\cite{atlasgnx}, exists only at \Kfull{} classes.

\begin{table*}[t]
\centering
\small
\setlength{\tabcolsep}{4pt}
\caption{Discriminants analyses build by summing the output scores of a pretrained tagger, against the number of classes in the label set (columns, finest first): Sophon's and the boosted $|V_{cb}|$ measurement's, and a discriminant that weights each QCD subclass separately, as GN3X evaluates its rejection per QCD subclass. $\bullet$: the discriminant can be built exactly from that label set's outputs, because a merged class's score is the sum of the scores of the classes it merges; ---: it cannot. This is a property of the label set, not of a trained model.}
\label{tab:usecase}
\begin{adjustbox}{max width=\linewidth}
\begin{tabular}{l l cccc}
\toprule
discriminant & source & 188 & 162 & 43 & 17 \\
\midrule
$X\to b\bar b$ vs QCD & \cite{sophon}, Eq.~(4) & $\bullet$ & $\bullet$ & --- & --- \\
$X\to c\bar s$ vs QCD (Sophon's $A_1$) & \cite{sophon}, Eqs.~(5)--(6) & $\bullet$ & $\bullet$ & --- & --- \\
three quarks, exactly one $b$, vs QCD (Sophon's $A_3$) & \cite{sophon}, Eqs.~(5)--(6) & $\bullet$ & $\bullet$ & --- & --- \\
$W'\to W\phi\to WWW$ event discriminant & \cite{sophon}, Eqs.~(7)--(9) & $\bullet$ & $\bullet$ & --- & --- \\
$D_{bc}$, boosted $|V_{cb}|$ measurement & \cite{vcbboosted}, Eq.~(1) & $\bullet$ & $\bullet$ & --- & --- \\
QCD subclasses weighted separately (GN3X-style) & \cite{atlasgnx}, Sec.~6 & $\bullet$ & --- & --- & --- \\
\bottomrule
\end{tabular}
\end{adjustbox}
\end{table*}


\label{sec:controls}% the generated caption of tables/appendix_vocabulary.tex refers to this label
Random partitions test whether, at a fixed number of classes and fixed loss weights, the cost of a coarse label
set depends on which classes are merged. A random partition is a label set with as many classes as the
\Kcoarse-class level, formed by assigning the resonant classes to groups at random; with one class for all QCD
jets the groups make \Kcoarse{} classes. Each group holds the share of the reweighted training jets of one
\Kcoarse-class group to within \RandVtwoShareTol, so it carries the same weight in the loss. A random partition
thus has the class count and the loss weights of the \Kcoarse-class level but not its physics
(figure~\ref{fig:design}(b)). The partitions are selected rather than typical draws. Each of \RandVtwoTrials{}
draws per partition is refined by a search that merges as few of the class pairs merged by the
\Kcoarse-class level as the share rule allows, and the draw with the fewest such pairs is kept. The partitions are therefore
built to cut across the \Kcoarse-class groups. They mix decays with different numbers of prongs and keep many
flavor boundaries.
Every conclusion we draw from them is restricted to partitions built this way. We pretrain
\RandVtwoNPartitions{} partitions, \RandVtwoRuns{} run each; appendix~\ref{sec:app-labels} gives their
construction and their groups.

Random partitions with controlled flavor boundaries test whether a single flavor boundary, at a fixed number of
classes and fixed group shares, accounts for the cost of the flavor merge. The three partitions differ only in
their flavor boundaries. In the random partition with no flavor boundary, every group is a union of whole
families of decays that differ only in quark flavor, with the group shares of the \Kcoarse-class groups. The
random partition with a $b$-quark boundary starts from it and moves into another group the
\RandFlavCutClasses{} of the \RandFlavSplitOrbitClasses{} four-quark decays \RandFlavSplitOrbit{} that contain a
$b$ quark. To keep the shares exact, \RandFlavOrbitClassesMoved{} decays of other forms move the other way in
whole families. Its one flavor boundary thus separates the four-quark decays with a $b$ quark from those
without. The random partition with a random flavor boundary makes the same compensating move but moves
\RandFlavCutClasses{} four-quark decays chosen at random, with group shares within \RandFlavFonerShareTol{} of
those of the $b$-quark partition. Each is pretrained with \RandFlavRuns{} seeds. Together with the random
partitions they are the controls; every control has \Kcoarse{} classes and is read against the
\Kcoarse-class level.

An auxiliary mass regression tests whether a physics target beyond the labels changes what the frozen features
carry, and what it costs the classification. The \Kres- and \Kcoarse-class models are also pretrained, with
\MlstMassNRunsWord{} seeds each, with the auxiliary mass regression. An extra output regresses
\begin{equation}\label{eq:masstarget}
  y = \ln\bigl(m_{\mathrm{gen}} / m_{\mathrm{SD}}\bigr),
\end{equation}
where $m_{\mathrm{gen}}$ is the soft-drop mass of the matched generator-level jet and $m_{\mathrm{SD}}$ that of
the reconstructed jet; only jets with a generator-level match enter the regression term. Its log-cosh loss is
added to the classification loss with the weight $\lambda = \MlstMassLambda$ at both label sets. This is the
value GloParT tuned for its two regression terms on generator-level masses~\cite[Eq.~(11.7)]{glopart}; the
target is ours, and $\lambda$ is not tuned on validation. The published CMS tagger also uses a log-cosh loss
but regresses the mass in GeV, with weight \LitCmsPartMassLambda~\cite{cmspart}, a value not comparable with
ours. At a fixed weight, the share of the training loss taken by the regression term differs between the
label sets because the classification loss is smaller with fewer classes. Averaged over training, it is
\MlstMassLossShareResMass\% in the \Kres-class runs and \MlstMassLossShareCoarseMass\% in the
\Kcoarse-class runs. The two label sets are therefore compared at a fixed weight, not at a matched share.

The \Kfull-class model with kinematic inputs only is a diagnostic of the input mismatch on top tagging, not a
method we propose. It is the \Kfull-class model trained on the particle four-momenta only, that is on the
\FillTopNPresent{} of the \FillNFeatures{} input features that the top-tagging data provide, as ParT
pretrains a model on kinematic features for top tagging~\cite{part}. Its pairwise features, its training jets
and their order, and its training settings are those of the other \Kfull-class models. Its initial weights are
its own, because its first layer is narrower. The kinematic-only model, pretrained once, answers two questions.
After fine-tuning, it measures how much top-tagging rejection the missing inputs cost, against the
\Kfull-class model fine-tuned with those inputs set to zero, which shares its rescaling, data and settings.
With frozen features and nothing missing, it shows whether a linear probe on top tagging is limited by the
frozen features or by the missing inputs.
% CHECK: pending the kinematic-only model (pretraining in progress; its linear probes and fine-tuning are not run yet)

% ----------------------------------------------------------------------------
\subsection{Evaluation}\label{sec:evaluation}

Four classification tasks each ask one question. JetClass-II at its \Kfull{} classes asks which distinctions
the frozen features keep. Its classes are the \Kfull-class model's own labels, so this task favors finer label
sets by construction; we also report it over the resonant and the QCD classes separately, since only the
\Kfull-class level divides the QCD jets. JetClass~\cite{part} has every input present but other classes and a
different simulation without pileup. It is the in-distribution downstream test and asks whether a label-set
effect survives on new classes and a new simulation when no input is missing. Top tagging~\cite{toptag} and
quark/gluon discrimination~\cite{efn,qgpythia} are the public benchmarks, with the companion Herwig
sample~\cite{herwigsevenone,qgherwig} as a change of generator. They lack some inputs, and they ask what the
label set changes where analyses compare taggers. The jet mass and the tagging of $X\to b\bar b$ and $X\to bs$
complete the evaluation. Dataset sizes, selections and splits are given with each result.

The network reads \FillNFeatures{} features per particle, and the public benchmarks lack some of them: the
top-tagging data provide only the four-momenta, from which \FillTopNPresent{} of the features follow, and the
quark/gluon data lack the track-displacement variables, ParT's ``exp'' scenario~\cite[Table~2]{part}
(sections~\ref{sec:top} and~\ref{sec:qg} list them). Features a benchmark does not provide are set to zero
for every label set, so top tagging and quark/gluon discrimination compare label sets on identical inputs, and
fine-tuning with these zeros is the benchmark comparison. Zero follows JetParticle-JEPA, which sets unused inputs to zero without retraining~\cite{jpjepa}; that model saw
zeros in pretraining through input dropout, and ours did not. OmniLearn argues against pretraining one model
for each dataset~\cite[Sec.~II.A]{omnilearn}. ParT instead pretrains for the inputs of each
benchmark~\cite{part}, which we follow only through the kinematic-only model.

A linear probe cannot adapt to an input that no pretraining jet contains, as a zeroed particle
is (section~\ref{sec:top}), so for linear probes on frozen features we compare three substitutions for the missing input features: zero, the training mean and random draws. The training mean is
the running mean of the network's input normalization, the centroid of the training particles. The random
draws are four draws of a real JetClass-II particle of a randomly drawn species, with the features averaged
over the draws. Zero is the controlled comparison. The training mean asks whether a value at the center of the
training distribution recovers what zero loses, and the random draws ask whether a value that a real particle
can take recovers it. One substitution per task is applied to every model, the randomly initialized network
included: the one whose linear probes have the lowest validation cross-entropy summed over all models. It is
\LpSubstitutionTop{} for top tagging and \LpSubstitutionQg{} for quark/gluon discrimination, whose choice the
Herwig jets share, and figure~\ref{fig:f5} shows the three substitutions side by side. % CHECK: pending LpSubstitutionTop, LpSubstitutionQg (substitution fits not pulled)

Linear probes are our main evaluation of what a label set puts into the frozen features, as linear evaluation
is for transferable features in vision~\cite{kornblith} and a frozen backbone with a linear head is for masked
particle modeling~\cite{mpm}. On JetClass-II the linear probes are trained and tested on class-balanced samples
from held-out files, with equal numbers of jets of each class where the data allow, so that rare classes
weigh as much as common ones. The training sets are nested; their sizes and the fit settings are in
appendix~\ref{sec:app-probes}. A linear probe reads only what a linear function of the features can. An MLP
probe, a small multilayer perceptron with one hidden layer of \LpMlpHidden{} units trained on the same frozen
features, therefore stands beside every linear probe that finds no difference, and tells whether the
information is absent or only not linearly readable. The randomly initialized network is the saved initial
weights of a \Kfull-class model before any training, one for each seed. Evaluated in the same way, it gives
the floor: what a linear probe on frozen features gets from the architecture alone. Its BatchNorm statistics
are recomputed on \DesignBnJets{} training jets, so that it is compared with the pretrained networks under the
same normalization of the inputs; appendix~\ref{sec:app-bn} repeats the pretrained networks' linear probes
with recomputed statistics.

Binary linear probes on pairs of JetClass-II classes ask whether a label set keeps a distinction it merges, or
one that a downstream discriminant needs. In two pairs, $X\to b\bar b$ against $X\to c\bar c$ and
$X\to YY\to bbqq$ against $X\to YY\to ccqq$, the decays differ only in quark flavor and are merged at
\Kcoarse{} classes. The third, $X\to bc$ against $X\to bq$ and $X\to cs$ together, sets the numerator class of
$D_{bc}$ against the resonant classes of its denominator~\cite[Eq.~(1)]{vcbboosted}; $X\to bc$ and $X\to bq$
share a class from \Kmid{} classes on. Each pair is merged in some random partitions and kept apart in others,
which separates the effect of merging a pair from that of the class count. The three pairs of QCD classes that
only the \Kfull-class level separates are in appendix~\ref{sec:app-pairs}, and the sizes of the pair samples
in appendix~\ref{sec:app-probes}.

Fine-tuning asks what label-set effect remains once every weight adapts. Each pretrained model is fine-tuned once, with its
whole head, hidden layer included, initialized anew. The same network trained on the downstream task alone,
from random initial weights, is the reference for what pretraining adds over the downstream data; it is
trained with \MlstFtRefSeedsWord{} seeds. On the full training sets of top tagging and quark/gluon
discrimination and on the largest JetClass training set, fine-tuning uses ParT's fine-tuning settings, and the
network trained from scratch uses ParT's from-scratch settings~\cite{part}. The smaller training sets use one
rule, identical for pretrained and from-scratch networks: a fixed maximum number of steps, early stopping on
the validation set and a learning rate chosen on validation (appendix~\ref{sec:app-ft}). Missing inputs are
set to zero in every fine-tuning. On JetClass-II the levels of the hierarchy and the models with the auxiliary
mass regression are fine-tuned, while the controls are evaluated with linear probes only. One fine-tuning per
pretraining run spends the compute on pretraining runs, whose variation this study must measure. Separating
the variation of pretraining from that of fine-tuning would need several fine-tuning runs per pretraining run,
as in MultiBERTs~\cite[App.~D]{multiberts}.

Each task is scored in the form its literature uses. On top tagging and quark/gluon discrimination the tables
give the accuracy, the AUC and the background
rejection $1/\epsilon_B$ at signal efficiencies of 50\% and 30\%, written $R_{50}$ and $R_{30}$, where
$\epsilon_B$ is the efficiency for the background. The text compares label sets by $R_{30}$, the background a
tagger removes at a working point an analysis uses, because values of the AUC close to one hide differences
that rejection shows. With pretraining, for example, ParT's top-tagging AUC rises only from
\LitPartTopAucScratch{} to \LitPartTopAucFt, while its $R_{30}$ rises from \LitPartTopRejScratch{} to
\LitPartTopRejFt~\cite{part}. On JetClass the tables follow ParT~\cite[Table~1]{part}. They give the accuracy,
the AUC averaged over all pairs of classes as in ParT, and the rejection of each signal class against the QCD
background at the signal efficiency ParT uses for it. The per-class rejection shows whether a label set helps
or hurts specific signals, in a form comparable with ParT. The AUC averaged over classes, each class against the rest,
is given beside it in appendix~\ref{sec:app-jetclass}. On JetClass-II the tables give the accuracy and the AUC
averaged over classes, each against the rest, over all \Kfull{} classes and over the resonant and the QCD
classes, the form in which the separation of many classes is usually reported.

Between pretraining configurations on JetClass-II, label sets are compared by the ratio of $1-\mathrm{AUC}$, the
fraction of signal--background pairs the classifier orders wrongly:
\begin{equation}\label{eq:ratio}
  \rho_{A/B} = \frac{1-\mathrm{AUC}_A}{1-\mathrm{AUC}_B}.
\end{equation}
Here $A$ and $B$ are two configurations scored on the same task and the same test jets, with the AUC averaged
over classes for the full task and the binary AUC for a pair of classes. A ratio above one means that $A$
orders more pairs wrongly than $B$. The ratio answers by what factor a configuration changes the errors on the
pretraining classes, where the AUCs lie close to one and differ by little in absolute terms. Rejection at a
fixed signal efficiency does not serve here, because for some pairs of classes the best models pass almost no
background jets of the test sample at the working point. The ratio of $1-\mathrm{AUC}$ appears only in this role: never
on JetClass, top tagging or quark/gluon discrimination, and never as a table headline.

No published jet foundation model reports results for every class it is pretrained on. Sophon and GloParT sum
classes into discriminants~\cite{sophon,glopartdp}, and OmniLearned evaluates its classes only through a summed
anomaly score~\cite{omnilearned}. We show all \Kfull{} JetClass-II classes in one heatmap of each class's
ratio of $1-\mathrm{AUC}$, the class against the rest, to the \Kfull-class model (figure~\ref{fig:f3}). The classes are
grouped by their \Kcoarse-class parent and the color scale is logarithmic, after the per-class heatmap of
JetParticle-JEPA~\cite{jpjepa}.

Three further descriptors each answer one question. The first is the relative drop in the accuracy of a
JetClass-II linear probe when the inputs a benchmark lacks are set to zero in its training and test jets. It
measures how much of what a linear probe reads depends on those inputs. JetParticle-JEPA reports the same
relative drop when it removes inputs at inference~\cite{jpjepa}; here the linear probe is also trained without
them. The second is the data-efficiency equivalence of a fine-tuned model: the number of training jets with
which the network trained from scratch reaches the same metric, found by linear interpolation in the logarithm
of the number of jets between neighboring training sizes. It asks how many from-scratch training jets
fine-tuning at a given size is worth. The third is the share of the gain from pretraining that a linear probe
keeps: its accuracy minus that of training from scratch, as a fraction of the accuracy of fine-tuning minus
that of training from scratch, at the same number of training jets. It asks which part of the gain survives
with the backbone frozen.

The jet-mass resolution asks how close a mass estimate $m_{\mathrm{est}}$ comes to the generator-level mass
$m_{\mathrm{gen}}$ of equation~(\ref{eq:masstarget}), at every jet $p_{\mathrm T}$ and mass. For the response
$r = m_{\mathrm{est}}/m_{\mathrm{gen}}$ we use the definition of ATLAS~\cite[Sec.~6]{atlasregression},
\begin{equation}\label{eq:resolution}
  \sigma_r = \frac{q_{+} - q_{-}}{2\,\tilde r},
\end{equation}
where $\tilde r$ is the median response and $q_{-}$ and $q_{+}$ bound the central 68\% of the distribution of
$r$; the median response is reported beside it. PELICAN uses the same quantile range of the relative residual
without the division by the median~\cite[Sec.~7.3]{pelican}, and CMS reports the position and the width of the mass
peak~\cite{cmspart,glopartdp}. The target $m_{\mathrm{gen}}$ is the same for resonant and QCD jets. We bin in
the reconstructed jet $p_{\mathrm T}$, where ATLAS bins in the true $p_{\mathrm T}$. The jets are JetClass-II
test jets with a generator-level match, divided into training, validation and evaluation parts for the
regressions. Three estimates are compared on the same evaluation jets: the soft-drop mass itself, the regressed
mass $m_{\mathrm{SD}}\,e^{\hat y}$, where $\hat y$ is the output of the auxiliary mass regression, and a
regression from frozen features. The last is an
MLP probe trained to predict $y$ after the training median of $y$ in each JetClass-II class is subtracted, so
that a better classifier is not credited as a better mass regression; it asks whether the frozen features
themselves carry the jet mass. A multilayer perceptron is used because the mass is not a linear function of
the features, and the linear (ridge) regression is in appendix~\ref{sec:app-mass}.

Mass sculpting asks whether selecting on the summed score of the resonant classes makes the background mass
distribution peak where a signal would. As CMS does~\cite[Sec.~6.2]{cmspart}, we select on
$1 - p_{\mathrm{QCD}}$, at QCD efficiencies of \SculptEffHigh{} and
\SculptEffLow. We then measure the Jensen--Shannon distance, the square root of the base-2 Jensen--Shannon
divergence, between the soft-drop mass distributions of the selected QCD test jets and of all QCD test jets;
figure~\ref{fig:f8b} also traces it over a grid of QCD efficiencies. A random selection of as many QCD jets
gives a distance above zero from the size of the sample alone. The expected distance of such a selection,
drawn without replacement, is therefore the floor against which each distance is read: it asks whether a
distance is larger than a random selection of the same size would give.

We repeat Sophon's resonance-tagging benchmark~\cite[Sec.~III]{sophon} at every label set, for a spin-0
resonance $X$ of mass \XbxMx~GeV decaying to $b\bar b$ or to $bs$. In the \Kfull- and \Kres-class models,
where $X\to b\bar b$ is its own class, it is tagged by Sophon's direct discriminant~\cite[Eq.~(4)]{sophon}. This
asks whether the pretrained output scores alone tag it without a new head. At every label set it is also
tagged by a transfer head: a multilayer perceptron on the frozen features whose first layer starts from the
pretrained head's hidden layer, trained with Sophon's settings~\cite[App.~B.5]{sophon}. $X\to bs$, a class in
no label set, is tagged by the transfer head and by the same network trained from scratch with Sophon's
settings for it~\cite[App.~B.3]{sophon}, with three seeds. The score of the output class that contains
$X\to bq$, the nearest pretrained class, is a control that asks whether a gain is in the new head or already
in the nearest output. The background is the QCD jets of the JetClass-II test files, with no trigger and no
jets from $W$ or $Z$ bosons or top quarks. Jets are counted in Sophon's soft-drop mass window,
$\LitSophonXWindowLo < m_{\mathrm{SD}} < \LitSophonXWindowHi$~GeV for a resonance of \LitSophonXMass~GeV,
scaled by the ratio of the resonance masses to \XbxWindowLo--\XbxWindowHi~GeV. We report
$\epsilon_S/\sqrt{\epsilon_B}$, which is proportional to the discovery significance for a small
signal~\cite{cowan} and so gives the reach of a selection without assuming event counts, and the rejection
$1/\epsilon_B$ at $\epsilon_S = 30\%$. Sophon reports the discovery significance at an integrated luminosity
of \LitSophonXLumi~fb$^{-1}$ on the full Standard Model background after a trigger, which needs the event counts
of every background process that our samples do not provide, so the two are not comparable in value. Each fit
records whether it uses a fixed or a variable resonance mass; Sophon trains its transfer head on variable masses and
evaluates it at \LitSophonXMass~GeV. % CHECK: pending XbxMx, XbxWindowLo, XbxWindowHi (X->bb/X->bs samples not staged; drop this paragraph with section 4.7 if not staged by 2026-11-12)

% ----------------------------------------------------------------------------
\subsection{Uncertainties}\label{sec:uncertainties}

Every uncertainty is the standard deviation over the independently pretrained networks of a configuration,
each followed by one linear probe or one fine-tuning. Each level of the class hierarchy has \MlstNRunsWord{}
pretraining runs. The models with the auxiliary mass regression have \MlstMassNRunsWord{} each, the random
partitions with controlled flavor boundaries \RandFlavRuns{} each, and each random partition and the
kinematic-only model one; the mean and standard deviation of the random partitions are taken over the
\RandVtwoNPartitions{} partitions. The spread therefore includes the variation between pretraining runs, which
most published foundation-model results leave out by repeating only the downstream training on one pretrained
network; Sophon reports no variation between pretraining runs~\cite{sophon}. Three runs are given as the mean
and standard deviation, rounded following the Particle Data Group~\cite{pdg}, two runs as both values and one
run as its value. With three runs a standard deviation is itself uncertain by about \MlstSdRelUnc. The network
trained from scratch varies its training seed, and the randomly initialized network its initial weights.

Ratios are formed between models with the same seed, which share their training jets and initial weights, and
then averaged over seeds; a control is paired with the models of the seed it was pretrained with. The task's
metric is $1-\mathrm{AUC}$ on JetClass-II, $R_{30}$ on top tagging and quark/gluon discrimination, and the rejection of
each signal class on JetClass. A paired ratio of it is a difference when it departs from one by more than twice
its standard deviation and by more than its test-set uncertainty, which a bootstrap of the test jets gives
(appendix~\ref{sec:app-boot}). JetClass rejections that rest on too few passing background jets, marked with a
dagger in table~\ref{tab:jetclass}, are not used for claims. For a configuration with two runs, a ratio is a
difference only when both runs lie further from one than twice the standard deviation of the same ratio between
levels of the class hierarchy. The finite test set adds less to the uncertainty than the variation between
pretraining runs (appendix~\ref{sec:app-boot}). % CHECK: pending MlstBootTestMaxProbe (class-balanced bootstrap) and MlstBootTestMaxFt; the test-set SD must lie below the seed SD of the JetClass-II label-set ratios
Training ran on several GPU types, and table~\ref{tab:compute} gives each run's types and GPU-hours. The \MlstNRunsWord{}
\Kfull-class models trained on one GPU type each, one of them on a different type from the other two. The
label-set ratios of that seed lie within \MlstGpuTypeBound{} of the mean of the other two seeds, no more than
the variation between seeds, so the GPU type does not move a result beyond that variation. % CHECK: MlstGpuTypeBound (today 1.4%) must not exceed the seed SD of the 17/188 ratio (MlstCmpJcTwoCoarseOverFull, today 1.9%); it is NOT far below the 162/188 effect (2.6%), so do not write "far below the label-set effects"
